% =============================================================================
%  «ТЕОРИЯ СИСТЕМ» — ствол серии «Архитектура Рассуждения»
%  Блок 8. ВЕТВЛЕНИЕ: теория становится мирами (ЗАКРЫТИЕ).  Разделы 8.1–8.5 (черновик прозы).
%  Стиль: «выводим и показываем»; репо-якоря — в приложении, не в сносках.
%  Самокомпилируемо: pdflatex blok-8-vetvlenie.tex (см. память latex-build-windows).
% =============================================================================
\documentclass[12pt,a4paper]{article}
\usepackage{cmap}
\usepackage[T2A]{fontenc}
\usepackage[utf8]{inputenc}
\usepackage[russian]{babel}
\usepackage{paratype}
\usepackage{amsmath,amssymb,amsthm}
\usepackage{listings}
\usepackage[expansion=false]{microtype}
\usepackage[margin=28mm]{geometry}

\newcommand{\rezultat}{\medskip\noindent\emph{Что отсюда следует.}\ }

\title{\textbf{Теория Систем}\\[2mm]\large Блок 8. Ветвление: теория становится мирами}
\author{}
\date{}

\begin{document}
\maketitle
\thispagestyle{empty}

\noindent\small\textit{Черновик: разделы 8.1--8.5 (закрытие тома). Продолжает Блок~7 (граница).
Теория замкнута на одной границе --- здесь показано, как из неё \emph{ветвятся} миры. Машинные
якоря --- в Приложении.}\normalsize

\bigskip

\noindent Теория собрана и замкнута на одной границе. Теперь --- куда она \emph{ветвится}. Каждая
ветвь есть E/R/R-разбор своей области, а H1 --- их сквозной нерв. И ветвление здесь не отсылка, а
\emph{связь}: мосты проведены машинно, ствол соединён с ветвями.

% =============================================================================
\section*{8.1\quad Математика --- главная ветвь}
% =============================================================================

Первая и главная ветвь --- математика. E/R/R разворачивается через все её области, а нерв всюду один:
\emph{финитизация $=$ конструктивность} (H1). Вещественное, алгебраическое замыкание, многообразие,
даже гипотеза Римана --- берутся как \emph{процессы}, а не как завершённые объекты.

\rezultat
Математика --- не приложение \emph{поверх} теории, а её \emph{разворачивание}: те же законы,
принципы и та же граница, развёрнутые по областям. Подробное --- в Томе «Математика».

% =============================================================================
\section*{8.2\quad Физика}
% =============================================================================

{\sloppy
Здесь мосты особенно конкретны. E/R/R-решёточное действие \emph{совпадает} с действием Вильсона
(тождество $2-2\cos\theta = \theta^2$); цепь E/R/R $\to$ Вильсон $\to$ характеры $\to$ трансфер
ведёт к десяткам наблюдаемых; различение $\to$ автоморфизм $\to$ генераторы смотрит в сторону
Стандартной модели.

Но честно и о зазорах. Это \emph{пере-выводы с зазорами}, а не новая физика; непрерывная группа
$SU(N)$ здесь \emph{постулируется} (нужна топология, которой в различении нет).\par}

\rezultat
Физика связана с ядром \emph{машинно}, но с \emph{honest-зазорами}: теория даёт язык и проверяемые
совпадения, не претендуя на новую физику.

% =============================================================================
\section*{8.3\quad Знание}
% =============================================================================

Знание разобрано как \emph{система} в E/R/R целиком: мост согласует базу знаний с формализацией,
а шесть выводов дают её костяк. Это --- образец того, как целая область укладывается в триаду.

\rezultat
Знание --- \textbf{образец} полного E/R/R-разбора области; подробное --- в Книге «Эпистемология».
Существующее познаваемо (§5.1) --- отсюда и мост к теории знания.

% =============================================================================
\section*{8.4\quad Рассуждение}
% =============================================================================

Заблуждения и парадоксы --- это \emph{нарушения} E/R/R: по горизонтали (смешение ролей) и по
вертикали (смешение правил и уровней). Не список ошибок, а сбои \emph{одной} структуры.

\rezultat
Ошибки рассуждения --- \textbf{нарушения одной структуры}, а не разрозненный каталог. Хорошая форма
(§2.9) читается тут как \emph{норма} рассуждения; подробное --- в Книге «Архитектура Размышления».

% =============================================================================
\section*{8.5\quad Потенциал и границы применения}
% =============================================================================

Скажем честно, чем теория \emph{является} и чем \emph{нет}. Она --- инструмент \emph{понимания и
верификации}, не ускорения. Для искусственного интеллекта она полезна на уровне \emph{рамки}:
онтология агентов и знаний, счёт ограниченной рациональности, таксономия эмерджентности, детектор
ошибок рассуждения.

\rezultat
Граница применения проведена честно: это \textbf{рамка понимания, а не движок}. Теория не
ML-алгоритм и не бенчмарк, и не претендует на новую математику или физику. Её сила --- в
\emph{ясности}, а не в скорости.

\medskip
\noindent Так замыкается том. \textbf{Книга открылась единицей} (§1.2) --- одним актом различения ---
\textbf{и закрылась границей и её мирами.} Обе ноты --- старт (позитивная единица) и предел
(role-limit) --- суть одна L4-асимметрия. Ствол отдал ветви; H1 держит их все.

% =============================================================================
\appendix
\section*{Приложение к Блоку 8 — машинные якоря}
% =============================================================================
\small\raggedright\sloppy
\noindent Всё перечисленное машинно проверено; единственная аксиома ветви --- \texttt{classic} (L3),
кроме честно отмеченного зазора $SU(N)$.

\begin{itemize}\itemsep2pt
\item \textbf{§8.1.} корпус Тома «Математика» (\texttt{ShrinkingIntervals\_ERR.v},
\texttt{CauchyReal.v}, \texttt{numbertheory/*}, \texttt{geometry/*}); H1-принцип
\texttt{FinitizationPrinciple.v}.

\item \textbf{§8.2.} \texttt{ERRWilsonBridge.v} (\texttt{two\_minus\_two\_cos},
\texttt{err\_equals\_wilson\_N1/N2}); \texttt{ERRComputationBridge.v} (\texttt{err\_is\_wilson});
\texttt{ERRGaugeFunctorSynthesis.v} (\texttt{gauge\_is\_automorphism}, \texttt{generators\_match\_sm};
$\triangle\,SU(N)$ постулируется).

\item \textbf{§8.3.} \texttt{ERRKnowledgeBase.v} (\texttt{err\_knowledge\_base\_full\_synthesis});
\texttt{foundation/Knowledge*.v}.

\item \textbf{§8.4.} \texttt{Architecture\_of\_Reasoning/}; хорошая форма
\texttt{ERRWellFormedness.v} (§2.9).

\item \textbf{§8.5.} потенциал и границы --- свод §13.
\end{itemize}
\normalsize

% --- Конец Блока 8. Том «Теория Систем»: черновая проза Блоков 1–8 завершена. ---

\end{document}
